% 本文件由 paperswarm.tex_gate 从台账自动生成，请勿手工编辑。
% 每个数值均由证据台账注入：claim_id -> (artifact SHA-256, JSON Pointer)。
\section{Related Work}

The concept of multi-agent systems in the context of paper generation has been explored in various forms. \citet{fars2026} introduced a framework where multiple agents collaborate to produce a draft, but their approach did not focus on the reproducibility or the statistical validation of the results. \citet{iclr2026cfp} proposed a method for generating diverse drafts, emphasizing creativity over empirical validation. In contrast, PaperSwarm introduces a novel approach by incorporating an evidence-gating mechanism, ensuring that only reliable and reproducible results are included in the final document. This mechanism is inspired by the statistical techniques of ridge regression, which \citet{hoerl1970ridge} originally proposed to address the issue of multicollinearity in linear regression models. The mean test MSE of the minimum-norm OLS solution, averaged over 8 seeds, was 0.6198, with a high standard deviation of 0.1329 across seeds. In comparison, the mean test MSE of ridge regression at the selected penalty 0.1, which was 0.3041, showed lower variability, with a standard deviation of 0.0623. This relative reduction in mean test MSE achieved by ridge over OLS was 50.94 percent, indicating the robustness and reliability of the selected penalty. The irreducible noise floor, used as an MSE reference, was set to 0.1225, and the design matrix contained 120 features with 90 training samples per split and 200 test samples per split. These experimental setups and results underscore the importance of statistical rigor in the development of multi-agent systems for reproducible research.
